\section{Check the last required counter before writing}
Counter exhaustion is a property of the whole request. For a positive byte count
\(n\), the number of blocks is \(b=\lceil n/64\rceil\), and the last required
counter is \(c+b-1\). The request is valid exactly when that value is at most
\(2^{32}-1\). A one-byte request at counter \texttt{ffffffff} is valid; a
65-byte request there is not. Rejecting every use of the largest counter would
unnecessarily discard its last block.

First enforce the local unsigned count limit. That makes \(n+63\) safe to compute
in a doubleword. Handle zero separately, then compute the last counter in 64 bits.
\texttt{lwu} zero-extends the starting counter, including values with bit 31 set.
A sign-extended \texttt{lw} used directly in this wide comparison would interpret
those counter bit patterns incorrectly.
\Needspace{13\baselineskip}
\MARSCodeLabel{Validate a positive bounded request before allocating its workspace.}
\begin{minted}[samepage=true]{mars-idea}
daddiu $t0, $a3, 63
dsrl $t0, $t0, 6
lwu $t1, 48($a0)
daddu $t0, $t0, $t1
daddiu $t0, $t0, -1
dsrl $t0, $t0, 32
bnez $t0, stream_counter_error
nop
\end{minted}

The shifted high part is zero when the last counter fits in one word. This check
precedes output stores, so a rejected request cannot leave a partly transformed
message behind. Checking for wrap only after writing the first block would violate
that interface.

\begin{table}[htbp]\centering
  \caption{Requests at the end of the counter range. Counts are decimal.}
  \tagpdfsetup{table/header-rows={1}}
  \begin{tabular}{lrrl}\toprule
    Starting counter  & Bytes & Blocks & Result                         \\\midrule
    \texttt{ffffffff} & 0     & 0      & Success without pointer access \\
    \texttt{ffffffff} & 64    & 1      & Success                        \\
    \texttt{ffffffff} & 65    & 2      & Reject, status two             \\
    \texttt{fffffffe} & 128   & 2      & Success                        \\
    \texttt{fffffffe} & 129   & 3      & Reject, status two             \\
    \bottomrule\end{tabular}\end{table}

The local state copy lets the transformation advance its counter without changing
the caller's state or nonce. Counter arithmetic never carries into word thirteen.
That word belongs to the nonce. Treating the last row as one 128-bit integer would
silently change which encryption the later blocks belong to.
